\originalchapter{辐射与物质及探测}\label{ch:detectors}
\originalsection{带电粒子的能损}
带电粒子穿过介质时与原子电子相互作用, 引发激发和电离. 对较重入射粒子, Bethe 公式描述中等到高能范围的平均碰撞阻止本领:
\begin{equation}
 -\frac{\dd E}{\dd x}=Kz^2\frac ZA\frac\rho{\beta^2}
 \left[\frac12\log\frac{2m_ec^2\beta^2\gamma^2T_{\max}}{I^2}-\beta^2-\frac\delta2\right].
\end{equation}
$z$ 是入射电荷数, $Z/A$ 是材料参数, $I$ 是平均激发能, $\delta$ 是密度修正. 此处 $K$ 的单位取决于 $A$ 的摩尔质量约定. $1/\beta^2$ 解释慢粒子的高电离能损, 对数上升解释相对论区的缓慢增加. 极低能需壳修正和有效电荷, 电子入射还需不同的运动学, 不应跨范围机械套用.

粒子末段能损增大形成 Bragg 峰. 平均射程近似 $R(E_0)=\int_0^{E_0}\dd E/S(E)$, 其中 $S=-\dd E/\dd x$; 能损涨落和多重散射使个体射程不完全相同. 核反应和辐射能损也会改变曲线.

\originalsection{光子的三种主要过程}
低到中能光子的光电吸收把全部光子能量交给束缚电子, 光电子动能约 $E_\gamma-E_b$; 吸收边与原子壳结合能相关. Compton 散射把一部分能量交给电子, 由能动量守恒得
\begin{equation}
 E'_\gamma=\frac{E_\gamma}{1+(E_\gamma/m_ec^2)(1-\cos\theta)}.
\end{equation}
电子能谱因此连续, 最大能损对应光子反散射. 当 $E_\gamma>2m_ec^2$ 时, 核场中可产生 $e^-e^+$ 对; 自由光子在真空中不能独自产生这对粒子, 因为无法同时守恒四动量, 核反冲承担必要动量.

不同过程的相对大小依能量和材料 $Z$ 变化: 光电作用在低能和高 $Z$ 更强, 中能常由 Compton 主导, 较高能对产生逐渐重要. 三条曲线的交点不是对所有材料都相同的常数. 不含材料和截面数据的图只能作为趋势示意.

\originalsection{轨迹、量热与粒子识别}
均匀磁场中曲率半径 $R$ 给横动量 $p_\perp=|q|BR$, 常用换算 $p_\perp[\mathrm{GeV}/c]\simeq0.2998\,|q/e|B[\mathrm T]R[\mathrm m]$. 磁场方向和弯曲方向一起给电荷符号. 有限测点分辨率与多重散射分别限制高动量和低动量区的精度; 单凭一条直轨迹不能区分电荷为零与动量极大.

量热器使粒子产生 shower 并测总沉积能. 电磁 shower 来自电子轫致辐射和光子对产生, 尺度用辐射长度 $X_0$; 强子 shower 含核碰撞与不可见能, 尺度用相互作用长度. 能量分辨率常经验写为 $\sigma_E/E=a/\sqrt{E}\oplus b\oplus c/E$, 各项分别来自统计涨落、校准非均匀和噪声, $\oplus$ 表示平方和后开方.

若速度满足 $\beta n>1$, 会产生 Cherenkov 辐射, $\cos\theta_C=1/(\beta n)$. 它不意味着粒子超过真空光速, 只超过介质中的相位速度. 将曲率动量与时间飞行或 Cherenkov 速度组合, 可由 $m=p\sqrt{1/\beta^2-1}/c$ 识别质量.

\originalsection{计数、效率与背景}
预期事例数为 $N=\mathcal L_{\rm int}\sigma\mathcal B\epsilon$, 积分亮度、产生截面、指定分支比和总效率都要出现. 若稀有衰变样本有 $N_p$ 个质子、观察时间 $t\ll\tau_p$, 则预期衰变数 $N_p t/\tau_p$. 零观察事件并不意味着寿命无穷: 零背景 Poisson 模型中 $P(0|\mu)=\e^{-\mu}$, $90\%$ 单侧上限约为 $\mu=2.30$, 对应寿命下限 $N_pt\epsilon/2.30$.

真实分析还需背景与效率误差; 本式仅给理想极限. 探测器设计、信号构造和计数统计把核与粒子的理论预测变成可检验的实验陈述.
